\documentclass{article}
\usepackage[utf8]{inputenc}
\usepackage{a4wide}

\begin{document}

\section*{เกมเดินป่าระหว่างอลิซกับบ๊อบ}

ในป่าลึกลับแห่งหนึ่ง มีต้นไม้โบราณที่มีโครงสร้างเป็นแบบ``ต้นไม้กราฟ'' (Tree

Graph)ต้นไม้นี้มีทั้งหมดnโหนด (จุด) โดยโหนดมีหมายเลขตั้งแต่0ถึงn --

1ต้นไม้มีรากอยู่ที่โหนด0และมีเส้นทางเชื่อมกันอยู่n - 1เส้น ซึ่งระบุในรูปแบบอาเรย์edges\\



แต่ละโหนดในต้นไม้จะมี``ประตูวิเศษ''ซึ่งจะต้องเปิดเมื่อมีคนเดินทางมาถึง\\

และคุณยังได้รับข้อมูลอีกชุดคือamountซึ่งบอกว่าแต่ละประตูมี``ราคา''หรือ``รางวัล''



\begin{itemize}

\tightlist

\item

  ถ้าamount{[}i{]}เป็น\textbf{ค่าลบ}คือ\textbf{ต้องจ่ายเงิน}เพื่อเปิดประตู

\item

  ถ้าamount{[}i{]}เป็น\textbf{ค่าบวก}คือ\textbf{จะได้รับเงินรางวัล}เมื่อเปิดประตู

\end{itemize}



\textbf{รูปแบบของเกม:}



\begin{itemize}

\tightlist

\item

  เริ่มต้น\textbf{อลิซ}อยู่ที่โหนด0และ\textbf{บ๊อบ}อยู่ที่โหนดbob

\item

  ทุก ๆ1วินาที:\\

\end{itemize}



\begin{itemize}

\tightlist

\item

  \begin{itemize}

  \tightlist

  \item

    \textbf{อลิซ}จะเดินไปยังโหนดถัดไปที่นำไปสู่\textbf{ใบไม้(leaf)}

  \item

    \textbf{บ๊อบ}จะเดินกลับเข้าหาโหนด\textbf{0}

  \end{itemize}

\item

  เมื่อใครมาถึงโหนดไหน:



  \begin{itemize}

  \tightlist

  \item

    ถ้าประตูยังไม่ถูกเปิด ต้องเปิด

  \item

    ถ้าถึงโหนดพร้อมกัน ต้อง``หารครึ่ง''ทั้งรางวัลหรือค่าจ่าย

  \item

    ถ้าประตูเคยถูกเปิดแล้ว ไม่ต้องจ่าย และไม่ได้เงินเพิ่ม

  \end{itemize}

\end{itemize}



\begin{itemize}

\tightlist

\item

  อลิซหยุดเดินทันทีเมื่อถึงโหนดปลายทาง (ใบไม้-leaf)

\end{itemize}



\begin{itemize}

\tightlist

\item

  บ๊อบหยุดทันทีเมื่อถึงโหนด0

\end{itemize}



\textbf{งานของคุณ}



หาว่า\textbf{อลิซจะมีรายได้สุทธิ (net income)}มากที่สุดเท่าไหร่



หากเธอเลือก\textbf{เส้นทางที่ดีที่สุด}ในการเดินไปยังโหนดปลายทาง (ใบไม้-leaf)



\textbf{ตัวอย่างที่1:}\\



\includegraphics[width=1.20833in,height=1.16667in,alt={image.png}]{/public/upload/6ddf9703e1.png}\\



\textbf{Input:}



edges = {[}{[}0,1{]},{[}1,2{]},{[}1,3{]},{[}3,4{]}{]}



bob =3



amount = {[}-2,4,2, -4,6{]}



\textbf{Output:}6



ลำดับเหตุการณ์:\\



\begin{itemize}

\tightlist

\item

  อลิซเริ่มที่โหนด0→เปิดประตู เสียเงิน -2→รายได้รวม = -2

\item

  บ๊อบเริ่มที่โหนด3

\item

  อลิซและบ๊อบเคลื่อนที่พร้อมกัน ไปยังโหนด1→ถึงพร้อมกัน→แบ่งรางวัล4คนละครึ่ง→ได้เพิ่ม

  +2→รายได้รวม =0

\item

  อลิซไปต่อที่โหนด3→ประตูถูกเปิดไปแล้วโดยบ๊อบ→ไม่ได้อะไร

\item

  บ๊อบมาถึงโหนด0→จบการเดินทางของบ๊อบ

\item

  อลิซไปต่อถึงโหนด4→เปิดประตู รับรางวัล +6→รายได้รวม =6

\end{itemize}



จบเกม: อลิซได้สุทธิ6ซึ่งเป็นค่าสูงสุด



\textbf{ตัวอย่างที่2:}



\textbf{\includegraphics[width=1.41667in,height=0.48958in,alt={image.png}]{/public/upload/cc0ab7ed9a.png}\\

}



\textbf{Input:}



edges = {[}{[}0,1{]}{]}



bob =1



amount = {[}-7280,2350{]}



\textbf{Output:}-7280



อลิซเดินจากโหนด0→1\\

บ๊อบเดินจากโหนด1→0



อลิซเปิดประตูที่โหนด0เสียเงิน -7280\\

ส่วนโหนด1ถูกบ๊อบเปิดก่อน→อลิซไม่ได้เงิน\\

รายได้สุทธิของอลิซคือ-7280



\textbf{ข้อกำหนด:}



\begin{itemize}

\tightlist

\item

  2 ≤ ~n ~≤ 100,000 ~ ~ จำนวน n คือจำนวน node ทั้งหมด

\item

  edgesมี n - 1รายการ

\item

  0 ≤ ai, bi \textless{} nและai ≠ bi

\item

  ต้นไม้ต้องมีโครงสร้างที่ถูกต้อง (ไม่มีวงจร)

\item

  1 ≤ bob \textless{} n

\item

  amount{[}i{]}เป็นเลขคู่ และอยู่ในช่วง{[}-10,000, 10,000{]}

\end{itemize}



\subsection*{Input}
บรรทัดที่ 1 : จำนวน i ตัวเชื่อมโหนด



บรรทัดที่ 2 - i+1 เส้นทางการเชื่อมต่อโหนด



บรรทัดที่ i+2 ค่าของ Bob



บรรทัดที่ i+3 : ค่าของรางวัล เรียงเป็นสายคั่นด้วยช่องว่างต้องแต่ node 0 - n



\subsection*{Output}
บรรทัดที่ 1 : ค่ารางวัลสูงสุดที่อลิซควรจะได้



\end{document}
